\section{Experiments: Part II}
\label{sec:exp2}

\subsection{Design: Beam value cannot be inferred from a Beam-only result}
\label{sec:exp2-design}

The Beam hypothesis (\cref{claim:beam}) is tested only by a matched three-way initialization
comparison at exact count and compute, all downstream settings identical:
\begin{itemize}[leftmargin=1.6em,itemsep=1pt]
  \item \textbf{Beam Fusion + carriers} (treatment, \cref{sec:beam-fusion,sec:carrier});
  \item \textbf{balanced compact carve} (strong image-free structural control,
        \cref{sec:exp1-baselines}) --- same input information, no tomographic fusion;
  \item \textbf{calibration-bounded random} (lower bound) --- uniform in the camera-derived
        bounding volume.
\end{itemize}
All three feed the identical compact optimizer and evaluation samples.  Reported:
initialization quality, fixed intermediate checkpoints, risk-curve area (time-to-quality),
selected endpoint, paired seed wins, and both development scenes plus
\toshow{fresh confirmatory scenes}.  Beam construction time and peak memory are accounted
separately --- the initializer must pay for itself end-to-end, not hide its cost in
``preprocessing''.  A converged tie after an initialization win means Beam improves startup
but not final capacity; the text will say exactly that if it happens.

\begin{table}[t]
  \centering\small
  \begin{tabular}{@{}lccccc@{}}
    \toprule
    Initialization & init $\Jq$ & AUC($\Jq$) & final $\Jq$ & held-out PSNR & seed wins \\
    \midrule
    Beam + carriers          & \tbd & \tbd & \tbd & \tbd & \tbd \\
    Balanced compact carve   & \tbd & \tbd & \tbd & \tbd & \tbd \\
    Bounded random           & \tbd & \tbd & \tbd & \tbd & \tbd \\
    \bottomrule
  \end{tabular}
  \caption{\redcaption{Matched three-way initialization comparison (per scene; paired seeds;
    Beam build time and peak memory reported in a companion column).  Entirely pending ---
    this is the claim-deciding table of Part~II.}}
  \label{tab:beam-value}
\end{table}

\begin{figure}[t]
  \centering
  \figplaceholder{5.5cm}{%
    \textbf{Figure 10 --- time-to-quality curves.}  Validation $\Jq$ (y, log) versus
    optimizer updates (x) for the three initializations of \cref{tab:beam-value}; seeds as
    thin lines, medians bold; annotate initialization points and the fixed checkpoints.  A
    second panel with held-out PSNR versus wall-clock (including Beam construction time as
    an x-offset for the Beam arm) makes the ``pays for itself'' question visible.  Data: the
    matched comparison runs.}
  \caption{\redcaption{Risk trajectories for Beam versus carve versus random initialization.
    To be produced.}}
  \label{fig:time-to-quality}
\end{figure}

\subsection{Ablations}
\label{sec:exp2-ablations}

\begin{table}[t]
  \centering\footnotesize
  \setlength{\tabcolsep}{4pt}
  \begin{tabular}{@{}lp{7.1cm}l@{}}
    \toprule
    Ablation & Development verdict (single scene) & Confirmatory \\
    \midrule
    No corrected covariance repair
      & repair $-63.2\%$ $\Jq$, 3/3 wins \evidence{C30}
      & \tbd \\
    Legacy (one-sided, undilated) repair
      & $+51.6\%$ $\Jq$: harmful, retired \evidence{C30}
      & --- \\
    Legacy opacity / appearance repair
      & $+9.8\%$ / $+0.12\%$: retired \evidence{C30}
      & --- \\
    One compact phase vs.\ two
      & second phase $\to 0.717\times$ $\Jq$, 3/3 \evidence{C30}
      & \tbd \\
    Phase-1 means frozen
      & fails $2\%$ non-inferiority \evidence{C30}
      & \tbd \\
    Means-only optimization
      & $4.18\times$ $\Jq$: rejected \evidence{C30}
      & --- \\
    Without containment prune
      & prune costs $2.7\%$ $\Jq$, zero violations \evidence{C30, C31}
      & \tbd \\
    Prune before vs.\ after phase 2
      & before selected \evidence{C30}
      & \tbd \\
    Phase-2 means trainable vs.\ frozen
      & frozen passes ($1.015\times$ $\Jq$) \evidence{C30}
      & \tbd \\
    Clone (incl.\ density-preserving)
      & fails $5\%$ gate: no growth \evidence{C30}
      & --- \\
    SH0 vs.\ incremental SH3
      & $2.3$--$3.0\% < 5\%$ floor: SH0 \evidence{C30}
      & \tbd \\
    Soft support (\cref{eq:soft-support})
      & hinge variant immaterial, retired
      & \toshow{freeze then test} \\
    \bottomrule
  \end{tabular}
  \caption{Carrier-stage ablation grid.  Development verdicts are frozen paired-root results
    on one outcome-exposed scene (19 fitting views, compact validation; evidence keys refer
    to \repopath{ara/logic/claims.md}).  \toshow{The confirmatory column --- fresh scenes,
    held-out cameras, $\ge 3$ paired seeds, preregistered gates --- is entirely pending;
    rejected legacy stages re-enter only as negative controls.}}
  \label{tab:ablations}
\end{table}

The development column of \cref{tab:ablations} closes the \emph{implementation-policy}
question on one scene: which stages the carrier pipeline runs, in which order.  It does not
close the \emph{value} question (\cref{tab:beam-value}), which is a comparison against
non-Beam arms, nor any cross-scene generalization.

\subsection{Diagnostics to report}
\label{sec:exp2-diagnostics}

Per Beam run: contributor-count histogram, per-view uniqueness, gate survival rates,
covariance eigenvalue spectra, lineage hashes; per maturation: carrier survival, per-family
optimizer motion, covariance residual distribution (\cref{eq:cov-residual}), containment
violation counts before/after prune, and the fixed-topology / frozen-mean containment
recheck.  These are the quantities that make a false-association failure (the designed risk
of CI fusion, \cref{sec:ci-fusion}) visible instead of silent.

\begin{figure}[t]
  \centering
  \figplaceholder{5cm}{%
    \textbf{Figure 11 --- Beam diagnostics panel.}  (a) contributor-count histogram; (b)
    carrier survival across stages (bar per stage); (c) distribution of covariance residuals
    \cref{eq:cov-residual} before/after repair; (d) containment violations per view before
    prune.  Data: diagnostics dictionaries already emitted by the implementation on the
    development scene (\repopath{runs/}).}
  \caption{\redcaption{Beam Fusion and carrier diagnostics.  To be produced from existing
    run artifacts.}}
  \label{fig:beam-diagnostics}
\end{figure}

\subsection{Decision rule for Part II}
\label{sec:exp2-decision}

Preregistered: Beam is claimed as a contribution only if it beats the balanced carve control
on the frozen primary endpoint (held-out quality at matched compute and capacity) with the
agreed seed-win rule, replicated on fresh scenes.  \toshow{The exact endpoint, materiality
margin, and seed-win rule must be frozen before the confirmatory runs (obligation O-9).}
A tie or loss is reported in full as a bounded result --- the tomographic construction,
observability result, and lineage machinery remain contributions of record, but not
quality claims --- and, per \cref{sec:decision-rule}, has no bearing on Part~I.
